\documentclass[12pt]{article}

\usepackage[T1]{fontenc}
\usepackage[spanish]{babel}
\usepackage[letterpaper, margin=2.54cm]{geometry}
\usepackage{csquotes}

\begin{document}

% Portada
\begin{titlepage}
	\centering
	{\scshape\LARGE Universidad Aztlán \par}
	\vspace{1cm}
	{\scshape\Large Campus Playa del Carmen \par}
	\vspace{3cm}
	{\scshape\Huge Cuentas principales del estado de resultados \par}
	\vspace{1cm}
	{\scshape\Large Resumen \par}
	\vspace{3cm}
	{\Large José Rosendo Camal Teh \par}
	\vfill
	{\Large Ingeniería en Sistemas Computacionales \par}
	\vspace{0.5cm}
	{\Large Segundo Cuatrimestre \par}
	\vspace{0.5cm}
	{\Large Contabilidad Financiera \par}
	\vfill
	{\Large Profa. Guadalupe Parada Valencia \par}
	\vfill
	{\Large 16 de abril de 2026 \par}
	\vfill
	{\Large Playa del Carmen, Quintana Roo \par}
\end{titlepage}

\section*{Resumen}

% Introducción
Las cuentas principales del Estado de Resultados son aquellas cuentas en las que se registran los montos de las transacciones de ingresos y gastos de una entidad. A continuación, se indagará de forma sintetizada en las siguientes: \textquote{Ventas Totales}, \textquote{Devoluciones sobre ventas}, \textquote{Descuentos sobre ventas}, \textquote{Compras}, \textquote{Gastos de compra}, \textquote{Devoluciones de compra}, \textquote{Descuentos sobre compras}, \textquote{Inventario inicial}, \textquote{Inventario final}, \textquote{Gastos de ventas o directos}, \textquote{Gastos de administración o indirectos}, \textquote{Gastos y productos financieros} y \textquote{Otros gastos y productos}.

Entre las diversas cuentas la primera mención se otorga a \textbf{Ventas Totales}. Esta tiene el objeto de contabilizar los valores ---en términos monetarios--- de las mercancías entregadas a los clientes de la entidad. Las mercancías son vendidas al contado o bien entregadas mediante crédito. Sin requerir una garantía documental para la transacción de venta que para algunos casos y, en otros, existe la necesidad del respaldo.

En desafortunadas ocasiones, los clientes devuelven las mercancías porque no les satisface su compra, sea por la calidad, el precio o por una particular característica, un criterio o un capricho. El valor de lo devuelto se registra en la cuenta de \textbf{Devoluciones sobre ventas}.

Así como hay desafortunados eventos, también hay casos favorables en los que una entidad otorga beneficios y rebajas a las mercancías. Existen distintas razones para que una entidad ejecute estas acciones ---de las que no indagaré---, pero todas comparten el beneficio en ambos actores: la entidad y su cliente. El monto de las bonificaciones y la cantidad de las rebajas se anotan en la cuenta de \textbf{Descuentos sobre ventas}.

Como se ha mencionado, los clientes pueden realizar adquisiciones a la entidad. De manera análoga, la entidad requiere ser un cliente para efectuar sus operaciones y para ello necesita comprar mercancías a otras entidades. Bajo esta perspectiva, a esas entidades se les conoce como \textquote{proveedores}. El importe de las compras se registra en la cuenta cuyo nombre nos es familiar: \textbf{Compras}.

De manera regular, las adquisiciones a los proveedores conllevan erogaciones adicionales: fletes y acarreos; derechos aduanales; seguros ---como la protección contra robos---; cargas y descargas, etc. Esta situación no es muy desconocida en la actualidad. Hoy en día, la actualidad se rodea de compradores y vendedores en un espacio virtual, sin límites geográficos ni limitaciones en cuanto a lo grande o pequeño de la mercadería. Estos montos se asientan en la cuenta \textbf{Gastos de compra} [vaya ironía de pagar para obtener lo que ya se ha comprado].

Lamentablemente, de manera similar como le puede ocurrir a los clientes, la entidad puede encontrarse en la situación de no estar satisfecha con las mercancías adquiridas, sea la razón que sea. Dado lo anterior, la entidad las devuelve a los proveedores y estas cantidades se anotan en la cuenta de \textbf{Devoluciones sobre compras}.

Hasta ahora, las cuentas mencionadas hacen referencia implícita a la actividad de poseer existencias de mercancías. No se puede vender ni adquirir mercancías sin que la entidad y los proveedores posean un \textquote{inventario}. En esta cuenta, que en realidad son dos, se registran los importes, los valores de las mercancías para ser vendidas y aquellas que tienen una finalidad distinta a la de la venta. Los montos correspondientes al inicio de un ciclo contable se contabilizan en \textbf{Inventario inicial}. Mientras que los valores obtenidos al concluir un periodo contable se anotan en \textbf{Inventario final}.

Las actividades de vender, comprar, inventariar y llevar la contabilidad en una entidad conllevan desembolsos de dinero como pago de sueldos y salarios, prestaciones de ley, pagos de alquiler o mantenimiento de inmuebles, permisos para realizar actividades económicas, entre otros. Estos gastos se dividen en dos grandes grupos: las erogaciones destinadas directamente a las tareas de vender, promocionar y desarrollar el volumen de venta; y las que de forma indirecta impactan en las primeras, como la administración y dirección de la entidad. 

El primero se asienta en la cuenta \textbf{Gastos de venta o directos} y el segundo grupo de gastos, en \textbf{Gastos de administración o indirectos}. La diferencia más notable es que, en la primera, los principales gastos en emolumentos van dirigidos al personal \textquote{operativo} y en la otra, al personal \textquote{administrativo} y \textquote{directivo} ---los trabajadores de cuello blanco---.

La penúltima cuenta es \textbf{Gastos y productos financieros}. Esta cuenta incluye los montos de gastos e ingresos usuales de la entidad. Ejemplos de ellos: intereses de diversos documentos, reducción de intereses por anticipos, intereses moratorios, intereses bonificados de créditos, pérdidas o utilidades en cambio de divisas extranjeras, comisiones bancarias, pagos por servicios bancarios y generales, etc.

Por último, la cuenta \textbf{Otros gastos y productos} ---a diferencia de la anterior--- contabiliza las operaciones no ordinarias de la entidad. Se enumeran algunos ejemplos: pérdida o utilidad en ventas de activos y en compraventas de acciones y valores; dividendos de acciones que posee la entidad y rentas cobradas.

% Conclusión
En total se explicaron las trece cuentas principales que, en conjunto, tienen el objetivo de registrar ---en cuanto a lo visto durante el curso--- las cantidades (en dinero) de las operaciones de ingresos y erogaciones de una entidad: ventas, devoluciones, descuentos, compras, entre otras.

\end{document}
